% Kelvin. Compare to a14 to isolate fillfactor.
\subsection{\expt{a11}: Fillfactor and HOT updates}
A11 retains a14's maintenance and sets \code{fillfactor = 80} on customers,
products, inventory and orders. Leaving free space can enable heap-only tuple
(HOT) updates when the relevant indexed columns are unchanged, avoiding new
entries in ordinary indexes \citep{pg-hot}. Kelvin collected update statistics
from worker shards before and after each measured run, outside the timed interval.

\begin{table}[htbp]
\caption{HOT share of updates across six client counts.}
\label{tab:a11}
\begin{tabularx}{\linewidth}{@{}YZZ@{}}
\toprule
\textbf{Table} & \textbf{a14: default fillfactor} & \textbf{a11: fillfactor 80} \\
\midrule
Inventory & 23.3--37.0\% & 100\% \\
Orders & 76.2--81.5\% & 100\% \\
Customers & 0--19.0\% & 100\% \\
Products & 1.8--12.8\% & 100\% \\
\bottomrule
\end{tabularx}
\end{table}

All recorded a11 updates were HOT and none moved to a new page. Nevertheless,
overall TPS and write latency did not improve consistently. Inventory received
only about 2,500--3,100 updates per ten-minute measurement across four workers,
because distributed reads still limited throughput. Avoiding index maintenance
at that rate had limited effect on total runtime.

Lower fillfactor also has a storage trade-off. T6's one-client median increased
from 266.425 to 280.749\,ms, consistent with a larger footprint, though this
experiment does not independently establish the cause. A11 recorded zero retries
at every client count; repetition is needed before attributing that observation
to fillfactor. Testing the change on a112 remains future work.
\evidence{\expt{a14}/job 967194 and \expt{a11}/job 967195:
\code{hot_updates.csv}, \code{transaction_metrics.csv} and \code{summary.csv}.}
